\documentclass[11pt]{article}

\usepackage[margin=24mm]{geometry}
\usepackage{amsmath,amssymb}
\usepackage{booktabs,longtable,tabularx,array}
\usepackage{enumitem}
\usepackage{xcolor}
\usepackage{hyperref}
\usepackage{microtype}
\usepackage[T1]{fontenc}
\usepackage{lmodern}

\hypersetup{
  colorlinks=true,
  linkcolor=blue!55!black,
  urlcolor=blue!55!black,
  citecolor=blue!55!black,
  pdftitle={PLUS and MOOR: Ricker-Assumption Stress Test in a Conservation Benchmark}
}

\setlist{leftmargin=*,itemsep=0.25em,topsep=0.35em}
\setlength{\parindent}{0pt}
\setlength{\parskip}{0.55em}
\setlength{\emergencystretch}{3em}
\renewcommand{\arraystretch}{1.16}

\newcommand{\method}[1]{\textsc{#1}}
\newcommand{\code}[1]{\texttt{\detokenize{#1}}}
\newcolumntype{L}[1]{>{\raggedright\arraybackslash}p{#1}}
\newcolumntype{Y}{>{\raggedright\arraybackslash}X}
\newcommand{\statusbox}[1]{%
  \begin{center}
  \fcolorbox{black!55}{yellow!10}{\parbox{0.92\textwidth}{#1}}
  \end{center}}
\newcommand{\takeaway}[1]{%
  \begin{center}
  \fcolorbox{blue!55!black}{blue!4}{\parbox{0.92\textwidth}{\textbf{Take-away.} #1}}
  \end{center}}

\title{PLUS and MOOR Under Hidden Demographics\\
\large From the Original Ideas to a Ricker-Assumption Stress Test in Real-Ecology Conservation Data}
\author{Supervisor-facing methods explanation and slide-preparation reference}
\date{20 July 2026}

\begin{document}
\maketitle

\begin{abstract}
This document explains, without requiring access to the code, the basic ideas behind PLUS and MOOR,
why the earlier benchmark implementations were appropriately described only as PLUS-inspired and
MOOR-inspired, and how the current implementations restore the defining algorithmic structures of
the papers. It then explains the conservation-specific adaptations and the current stress-test
design: both ecological methods assume Ricker dynamics, Ricker environments form the in-family
control, and Allee, theta-logistic, and regime-switching environments form three genuine
out-of-family tests. The final sections separate implementation validity from empirical performance
and identify concrete improvements for a later, stronger study.
\end{abstract}

\statusbox{\textbf{Current evidence status.} The scientific methods described here are
paper-aligned adaptations, not exact reproductions. The earlier cross-family PLUS configuration is
a separate frozen extension. The current Ricker-only PLUS production design contains 72 cells and
eight Ricker candidates per cell. Its first reduced structural smoke test stopped before fitting
because a launch parser attempted to convert the truthful sentinel
\code{regime_path_count=not_applicable} to an integer. No production job or performance inspection
occurred. The parser must be fixed, provenance refrozen where affected, and the smoke test passed
before the production experiment is described as executed or accepted.}

\begingroup
\small
\setlength{\parskip}{0pt}
\tableofcontents
\endgroup
\newpage

\section{The research question in one page}

The benchmark asks how to choose conservation interventions when important demographic quantities
are hidden. The manager sees historical action--survey episodes and receives future noisy surveys,
but does not receive the true ecological family, true demographic parameters, or the private safety
state. The scientific comparison is between two broad sources of strength:

\begin{itemize}
  \item \textbf{Ecological inductive bias:} use mechanistic population models to constrain what
  dynamics are considered plausible.
  \item \textbf{Black-box flexibility:} learn transitions, values, or policies with fewer explicit
  ecological commitments.
\end{itemize}

PLUS and MOOR are useful ecological comparators because they use mechanistic knowledge differently.
PLUS retains several fixed mechanistic explanations and updates their probabilities online. MOOR
fits one mechanistic explanation offline and commits to it. This distinction remains meaningful even
when both methods assume the same Ricker family.

The present stress test deliberately gives both methods the same family-level assumption:

\begin{center}
\begin{tabular}{lll}
\toprule
\textbf{True hidden environment} & \textbf{Ricker-only PLUS} & \textbf{Ricker MOOR} \\
\midrule
Ricker & In-family control & In-family control \\
Allee & Out-of-family misspecification & Out-of-family misspecification \\
Theta-logistic & Out-of-family misspecification & Out-of-family misspecification \\
Regime-switching & Out-of-family misspecification & Out-of-family misspecification \\
\bottomrule
\end{tabular}
\end{center}

\takeaway{There are four hidden environment families in total: one assumed family (Ricker) and
three non-Ricker stress-test families. The family label is hidden in every cell, including the
Ricker control.}

\section{Minimal POMDP reminder}

A partially observable Markov decision process (POMDP) separates the unobserved ecological state
from the noisy observation available to the manager. At time $t$:

\begin{enumerate}
  \item the system has a latent state $s_t$, including population abundance and relevant context;
  \item the manager holds a belief $b_t(s)=P(s_t=s\mid\text{history})$;
  \item the manager chooses conservation action $a_t$;
  \item the ecosystem transitions under $P(s_{t+1}\mid s_t,a_t)$;
  \item the manager observes survey $o_{t+1}$ under $P(o_{t+1}\mid s_{t+1})$; and
  \item Bayes' rule updates the belief before the next decision.
\end{enumerate}

A POMDP policy must account for both ecological consequences and information. An action can be
valuable because it improves the population, because it produces informative observations, or both.
This is precisely where a belief-state planner differs from a fully observed controller or a QMDP
shortcut that largely assumes uncertainty will disappear after one step.

\section{Published PLUS: the basic idea}

\subsection{Uncertainty over complete mechanistic models}

PLUS represents scientific uncertainty with a finite bank of complete candidate POMDPs
$M_1,\ldots,M_J$. A candidate is not merely a regression curve or a label. It contains a complete
transition law, observation law, reward interface, discretized state representation, and fixed
ecological parameterization. Each candidate is solved separately.

PLUS maintains two types of uncertainty:

\begin{itemize}
  \item a within-candidate state belief $b_t^j$ for each model $M_j$; and
  \item a model posterior $w_t^j=P(M_j\mid\text{history})$ over candidates.
\end{itemize}

After action $a_t$ and observation $o_{t+1}$, candidate $j$ supplies its own predictive evidence:

\begin{equation}
w_{t+1}^j \propto w_t^j
p(o_{t+1}\mid b_t^j,a_t,M_j),
\qquad \sum_{j=1}^J w_{t+1}^j=1.
\end{equation}

The action is selected from posterior-weighted candidate values:

\begin{equation}
a_t^*=\arg\max_a\sum_{j=1}^J w_t^j Q_j(b_t^j,a).
\end{equation}

Parameters remain fixed during deployment. ``Learning'' in PLUS is sequential Bayesian
discrimination among the registered candidates, not continuous online refitting.

\subsection{What the paper requires---and what it does not}

The defining PLUS structure requires complete fixed candidate POMDPs, separate solutions, separate
beliefs, sequential model evidence, and posterior-weighted decisions. It does \emph{not} require
the bank to contain several ecological families, and it does not require exactly 16 candidates.
A finite bank of distinct Ricker parameterizations is therefore still recognizably PLUS.

The original work used SARSOP-based policy representations in a fisheries setting. The present
implementation uses finite-horizon PBVI in a conservation setting. That solver substitution must be
disclosed and checked, but it does not replace the Bayesian candidate-model structure.

\section{Published MOOR: the basic idea}

MOOR makes a different commitment. It fits one mechanistic population model offline from an ordered
historical trajectory, integrates latent process uncertainty during fitting, builds a POMDP from the
selected fit, and then plans in that same fitted model.

The key invariants are:

\begin{enumerate}
  \item retain the order of the historical sequence rather than treating rows as exchangeable;
  \item propagate latent biomass or abundance through the complete trajectory;
  \item estimate a mechanistic population model under process uncertainty;
  \item use normalized, bounded, multi-start optimization; and
  \item use the selected model consistently for kernel construction, filtering, and planning.
\end{enumerate}

MOOR does not maintain an online posterior over competing families. In the current benchmark it is
deliberately restricted to one Ricker model in every environment cell. The method is therefore
correctly specified at the family level in Ricker cells and deliberately misspecified elsewhere.

\section{Why the previous baselines are called ``inspired''}

The earlier implementations were useful, computationally convenient exploratory baselines. Their
results remain valid for the algorithms actually executed. The issue is naming: their internal
objects were easier approximations of the papers' ideas rather than the papers' defining inference
structures.

\subsection{Previous PLUS-inspired baseline}

The earlier PLUS-inspired method fitted action-conditional polynomial transition templates labelled
Ricker, Allee, theta, and regime. Those labels did not instantiate the corresponding ecological
equations or complete candidate POMDPs. The method approximately reweighted templates and used a
tabular/QMDP-like planning mechanism.

It retained the broad motif ``maintain several explanations and reweight them,'' but not the full
PLUS mechanism of candidate-specific mechanistic likelihoods, beliefs, and separately solved
belief-state policies. The accurate description is therefore:

\begin{quote}
PLUS-inspired posterior model averaging over action-conditional polynomial transition templates.
\end{quote}

\subsection{Previous MOOR-inspired baseline}

The earlier MOOR-inspired method used independent per-action log-linear ridge regressions, for
example

\begin{equation}
\log(1+o_{t+1}/S)=\beta_{a0}+\beta_{a1}\log(1+o_t/S).
\end{equation}

It treated transitions as exchangeable one-step pairs and then converted the regressions into a
tabular transition model. It did not fit one mechanistic growth/capacity model through complete
ordered latent trajectories. The accurate description is therefore:

\begin{quote}
MOOR-inspired per-action regression baseline with tabular model-based control.
\end{quote}

\begin{longtable}{L{0.20\textwidth}L{0.35\textwidth}L{0.35\textwidth}}
\toprule
\textbf{Dimension} & \textbf{Previous inspired implementations} & \textbf{Current paper-aligned adaptations} \\
\midrule
\endfirsthead
\toprule
\textbf{Dimension} & \textbf{Previous inspired implementations} & \textbf{Current paper-aligned adaptations} \\
\midrule
\endhead
Ecological model & Generic fitted regressions/templates & Registered mechanistic Ricker equations and parameters \\
Historical data & Exchangeable or mostly one-step rows & Ordered complete episodes with explicit resets \\
PLUS uncertainty & Approximate template weights & One belief and solved value representation per fixed candidate, plus a Bayesian model posterior \\
MOOR commitment & Per-action regressions & One jointly fitted Ricker trajectory model used throughout \\
Planning & Tabular/QMDP-like approximation & Belief-state finite-horizon PBVI \\
Defensible label & PLUS-inspired or MOOR-inspired & Paper-aligned PLUS or MOOR adaptation \\
\bottomrule
\end{longtable}

\section{The conservation benchmark}

\subsection{State normalization and observations}

Let $S$ be a public survey-derived scale. Private abundance $N_t$, capacity $K_t$, and observed
survey $o_t$ are normalized as

\begin{equation}
x_t=N_t/S,\qquad k_t=K_t/S,\qquad y_t=o_t/S.
\end{equation}

For action $a$, let $u_a$ denote translocation/stocking, $d_a$ denote a capacity increment, and
$r_a$ denote a signed action-rate effect. The registered transform separates positive growth and
negative removal effects:

\begin{equation}
g_a=\max(r_a,0),\qquad h_a=\max(-r_a,0).
\end{equation}

Capacity and post-intervention abundance evolve as

\begin{equation}
k_{t+1}=\operatorname{clip}(k_t+d_a,k_0,k_{\max}),
\qquad m_t=\max(x_t+u_a,0).
\end{equation}

The survey model is lognormal:

\begin{equation}
y_t\mid x_t\sim\operatorname{LogNormal}(\log x_t,\sigma_o^2).
\end{equation}

The manager observes the noisy survey, not the true abundance or hidden family.

\subsection{Four hidden environment families}

The assumed Ricker map is

\begin{equation}
x_{t+1}=m_t\exp\left\{g_a\left(1-\frac{m_t}{k_{t+1}}\right)-h_a\right\}.
\end{equation}

The Allee environment introduces a low-abundance threshold $C$:

\begin{equation}
x_{t+1}=m_t\exp\left\{g_a\left(1-\frac{m_t}{k_{t+1}}\right)
\left(\frac{m_t}{C}-1\right)-h_a\right\}.
\end{equation}

The theta-logistic environment changes the density-dependence shape:

\begin{equation}
x_{t+1}=\left[m_t+g_am_t\left(1-
\left(\frac{m_t}{k_{t+1}}\right)^\theta\right)\right]_+\exp(-h_a).
\end{equation}

The regime-switching environment includes an unobserved discrete regime $z_t$ with transition
matrix $\Pi$; the active regime changes the growth mechanism or parameters. The current Ricker-only
PLUS and MOOR models contain neither $z_t$ nor $\Pi$. This is deliberate misspecification rather
than an implementation omission.

\section{How the papers are adapted to this benchmark}

\begin{longtable}{L{0.23\textwidth}L{0.29\textwidth}L{0.38\textwidth}}
\toprule
\textbf{Paper setting} & \textbf{Conservation adaptation} & \textbf{Reason} \\
\midrule
\endfirsthead
\toprule
\textbf{Paper setting} & \textbf{Conservation adaptation} & \textbf{Reason} \\
\midrule
\endhead
Fishing effort & Eleven conservation interventions & Actual decisions alter growth, mortality, habitat capacity, or abundance. \\
Catch observation & Noisy abundance survey & Catch is not observed in this benchmark. \\
Harvest/catch reward & Common conservation reward & All methods must optimize the same persistence, abundance, cost, and safety objective. \\
One historical series & Multiple ordered 25-step episodes & The offline data contain genuine reset boundaries. \\
Fishery biomass scale & Public survey-derived scale $S$ & Private demographic scale cannot be leaked. \\
Stationary context & Capacity, previous survey, and timestep & Habitat effects accumulate and the decision horizon is finite. \\
SARSOP/DESPOT & Context-conditioned finite-horizon PBVI & This is the available disclosed planner substitution. \\
\bottomrule
\end{longtable}

These adaptations change the application interface, but they do not replace the algorithms'
identities. The honest claim is ``paper-aligned conservation adaptation,'' not ``exact
reproduction.''

\section{Current paper-aligned implementations}

\subsection{Ricker-only PLUS}

The current method identifier is \code{plus_adapted_ricker_only_pbvi}. For each environment cell it
constructs eight complete fixed Ricker candidate POMDPs:

\begin{itemize}
  \item candidate 0: deterministic MAP fit using the full ordered fitting history;
  \item candidates 1--7: deterministic episode-bootstrap MAP fits;
  \item initial model prior: $w_0^j=1/8$; and
  \item online parameters: fixed; only beliefs and candidate probabilities update.
\end{itemize}

Every candidate receives its own transition and observation kernels, PBVI solution, and state
belief. Allee thresholds, theta exponents, regime multipliers, regime states, and $\Pi$ are
inapplicable. They must not affect any Ricker-only fit, kernel, filter, PBVI value, or action.

This bank represents \emph{parameter uncertainty within Ricker}, not uncertainty over the true
family. It therefore creates the intended family-misspecification stress test.

\subsection{Ricker MOOR}

The method identifier is \code{moor_adapted_ricker_misspec_pbvi}. It fits one Ricker model jointly
to complete training episodes. Each episode resets initial abundance and cumulative capacity, while
order within the episode is preserved.

Fitting uses process-only Monte Carlo. Because the observation is lognormal, its conditional mean is

\begin{equation}
\widehat y_{e,t}=E[y_{e,t}\mid x_{e,t}]
=x_{e,t}\exp(\sigma_o^2/2).
\end{equation}

The fitted objective is schematically

\begin{equation}
L(\vartheta)=\frac{1}{|\mathcal E|}
\sum_{e\in\mathcal E}\frac{1}{T_e}
\sum_{t=1}^{T_e}
E_{\text{process paths}}
\left[(\widehat y_{e,t}(\vartheta)-y_{e,t})^2\right].
\end{equation}

No independent observation-noise draw is inserted into this squared error. The minimum finite
training-objective fit across deterministic starts is used unchanged for kernel construction,
filtering, and PBVI.

\subsection{Same family assumption, different inference}

\begin{center}
\begin{tabularx}{\textwidth}{L{0.24\textwidth}YY}
\toprule
 & \textbf{Ricker-only PLUS} & \textbf{Ricker MOOR} \\
\midrule
Offline commitment & Several fixed Ricker candidates & One selected Ricker fit \\
Online model learning & Bayesian posterior over candidates & None \\
State uncertainty & One belief per candidate & One belief in the fitted model \\
Decision & Posterior-weighted candidate PBVI values & PBVI values from the single fitted model \\
Misspecification & All candidates wrong outside Ricker & Single model wrong outside Ricker \\
\bottomrule
\end{tabularx}
\end{center}

\section{Why the current implementation is truthful to the papers}

``Truthful'' here means that the defining algorithmic objects are present and consistently used;
it does not mean that the implementation copies every variable, solver, or hyperparameter from the
fishery studies.

\begin{longtable}{L{0.29\textwidth}L{0.28\textwidth}L{0.33\textwidth}}
\toprule
\textbf{Paper invariant} & \textbf{Current evidence} & \textbf{Disclosed adaptation} \\
\midrule
\endfirsthead
\toprule
\textbf{Paper invariant} & \textbf{Current evidence} & \textbf{Disclosed adaptation} \\
\midrule
\endhead
Complete PLUS candidate POMDPs & Eight serialized Ricker candidates with candidate-specific kernels & Ricker-only bank rather than the earlier cross-family extension \\
Separate PLUS solutions & Every candidate is independently solved and filtered & PBVI replaces SARSOP \\
Sequential PLUS evidence & Normalized candidate posterior updated from action--survey likelihoods & Surveys replace catches \\
Fixed PLUS parameters online & Candidate parameters frozen before deployment & Bootstrap MAP constructs candidate diversity \\
Ordered MOOR fitting & Complete episodes retain order and reset state/context & Multiple episodes replace one fishery series \\
One MOOR model throughout & Selected Ricker fit reused consistently in kernel, filter, and planner & PBVI replaces DESPOT \\
Process uncertainty & Monte Carlo latent propagation in fitting & Finite registered path budget \\
Privacy & Method context excludes hidden family and private demographics & Public action-category structure is retained \\
\bottomrule
\end{longtable}

The implementation also passed a small exact-POMDP probe in which PBVI chose the exact optimal
information-gathering action and matched that action's exact value within $10^{-5}$. This supports
the planner's ability to represent information value on that probe. It does not prove universal
PBVI exactness.

\section{Current 72-cell Ricker-assumption stress test}

\subsection{Registered design}

Only PLUS requires new execution for this deadline experiment. The balanced design is

\begin{equation}
9\ \text{populations}\times4\ \text{hidden families}\times
2\ \text{noise levels}=72\ \text{cells}.
\end{equation}

The selected survey-noise levels are $\sigma_o\in\{0.1,0.2\}$. Population coverage has priority
over adding the extreme noise levels. Each cell uses the safe reward configuration, the inherited
4,000-transition data budget, and an episode-level 80/20 fit/holdout split.

The registered computational settings include:

\begin{itemize}
  \item eight deterministic Ricker candidates and uniform prior $1/8$;
  \item eight L-BFGS starts and at most 100 iterations for production fitting;
  \item 16 process paths;
  \item 41 abundance bins, nine capacity bins, and 41 observation bins;
  \item 256 transition and 256 observation samples;
  \item PBVI horizon five, 32 belief points, seven observation branches, and discount 0.95; and
  \item one CPU per cell worker, with no candidate-level CPU sharding.
\end{itemize}

The 72 workers may run concurrently, but each worker processes its candidates sequentially and then
plans only if its fit succeeds. This caps PLUS at 72 CPUs, leaving cluster capacity for the separate
general-RL reruns.

\subsection{What this experiment can and cannot show}

If structurally accepted, the experiment can compare in-family versus out-of-family behavior across
nine populations at two moderate survey-noise levels. It can test whether online discrimination
among Ricker parameterizations gives PLUS robustness that differs from MOOR's one-model commitment.

It cannot by itself establish:

\begin{itemize}
  \item robustness at zero noise or the highest registered noise level;
  \item adequacy of eight candidates relative to larger banks;
  \item adequacy of 4,000 transitions;
  \item equivalence to the full 288-cell benchmark; or
  \item superiority of ecological methods over general RL before matched reruns and analysis.
\end{itemize}

Existing MOOR artifacts should not be rerun merely to create a new directory identity. They may be
paired later only where method, data, split, reward, PBVI, evaluation, and artifact hashes match and
method-specific return-blind acceptance passes.

\section{Implementation verification versus empirical evidence}

The distinction is essential for presentation:

\begin{itemize}
  \item \textbf{Implementation verification} asks whether the intended equations, privacy rules,
  beliefs, posterior updates, planners, hashes, and routing are correct.
  \item \textbf{Structural run acceptance} asks whether every registered cell completed under the
  frozen code/manifests and emitted valid return-blind diagnostics.
  \item \textbf{Performance analysis} asks how the accepted policies actually performed. It begins
  only after acceptance and explicit authorization to open returns.
\end{itemize}

The Ricker-only route added tests for candidate cardinality, bootstrap construction, privacy,
non-Ricker-constant invariance, cache identity, route separation, dispatch, and acceptance. However,
the first structural smoke test exposed a manifest parsing defect before any fit began:
\code{regime_path_count=not_applicable} was incorrectly passed to a generic integer conversion.
This is an operational/configuration routing defect rather than evidence against the PLUS
algorithm. A type-aware correction and passing smoke rerun are required before production.

No performance claim should be inferred from a test count, a successful smoke test, or a completed
fit alone.

\section{How to interpret eventual results}

\begin{itemize}
  \item Strong Ricker-cell performance but weaker non-Ricker performance would be evidence of
  sensitivity to family misspecification under the registered adaptation.
  \item Weakness in all families could arise from poor action-effect identification, observation
  noise, limited data, an inadequate candidate bank, or PBVI/discretization limits; it would not
  isolate one cause.
  \item PLUS outperforming MOOR outside Ricker could reflect useful parameter-model averaging and
  online evidence, even though every PLUS candidate is family-misspecified.
  \item MOOR outperforming PLUS could reflect the stability of committing to one well-fitted model
  rather than posterior dilution across weak candidates.
  \item Safe-mode weakness can be information-limited because the private safety objective is not
  supplied to the learner. It must not automatically be called intrinsic algorithm failure.
\end{itemize}

Results from the earlier cross-family PLUS bank answer a different question: that bank already
contains the correct family, so it tests structural-family discrimination rather than robustness to
a shared Ricker assumption. The two configurations require explicit version labels and must not be
silently pooled.

\section{What should be improved next}

\subsection{Data adequacy}

The 4,000-transition budget was inherited for comparison with general learners, not because it was
shown sufficient. It corresponds to 160 complete 25-step episodes, with approximately 128 episodes
used for fitting. Some action effects, especially capacity actions in sink populations, have much
weaker coverage. A future 4,000-versus-8,000 complete-episode sensitivity is necessary.

\subsection{Candidate-bank adequacy}

Eight Ricker candidates are a registered practical choice, not a demonstrated optimum. Future
work should examine nested bank sizes, bootstrap diversity, effective parameter-space coverage,
boundary concentration, and whether several candidates collapse to near-identical dynamics.

\subsection{Fitting diagnostics}

Future runtime instrumentation should emit per-start optimizer termination, projected gradients,
boundary-active coordinates, train/holdout trajectory errors, parameter uncertainty summaries, and
clear identifiability diagnostics. These cannot be reconstructed honestly from performance returns.

\subsection{Planner adequacy}

PBVI should be tested over belief-point count, observation branching, horizon, abundance/capacity
grids, and Monte Carlo kernel sample sizes. The exact toy probe is reassuring but does not establish
adequacy in every ecological cell.

\subsection{Model-class extensions}

After the primary Ricker-assumption stress test, useful named sensitivities include:

\begin{itemize}
  \item comparing Ricker-only PLUS with the separate cross-family PLUS extension;
  \item richer within-Ricker parameter grids or hierarchical priors;
  \item learned regime transition matrices using particle or stochastic EM in a separately named
  cross-family method;
  \item posterior predictive checks and calibrated candidate expansion; and
  \item a model-robust or nonparametric ecological extension that does not assume the true family is
  present.
\end{itemize}

\subsection{Experimental coverage and comparators}

Later work should restore $\sigma_o=0$ and $0.4$, complete the full registered cell matrix, rerun
general learners at exactly matched data/evaluation budgets, and report uncertainty across data and
evaluation seeds. Known-demographics work should distinguish known capacity from a stronger
known-action-effects oracle; action-specific $r_a$ values must not be described as one intrinsic
population growth parameter.

\section{Suggested supervisor-presentation sequence}

The following sequence can be converted directly into slides:

\begin{enumerate}
  \item \textbf{Question:} Can ecological structure help offline conservation control when
  demographics and the true model family are hidden?
  \item \textbf{POMDP reminder:} latent abundance, noisy surveys, belief updates, and actions that
  affect both ecology and information.
  \item \textbf{PLUS idea:} several fixed mechanistic POMDPs, one belief per candidate, online
  posterior, posterior-weighted decisions.
  \item \textbf{MOOR idea:} one mechanistic model fitted to ordered trajectories, then used
  consistently for belief planning.
  \item \textbf{Why the old versions are ``inspired'':} simpler polynomial/ridge models retained
  broad motifs but not the defining paper objects.
  \item \textbf{Conservation adaptation:} interventions, abundance surveys, episode resets,
  normalized states, common reward, and PBVI.
  \item \textbf{Fair stress test:} both methods assume Ricker; show the one-control/three-stress
  family table.
  \item \textbf{Current experiment:} eight Ricker candidates, nine populations, four families, two
  noise levels, 72 PLUS cells, matched existing MOOR rows where provenance permits.
  \item \textbf{Why it is truthful:} invariant-to-implementation table and verification/acceptance
  separation.
  \item \textbf{Limits and next work:} data adequacy, candidate diversity, optimizer diagnostics,
  PBVI sensitivity, extra noise levels, and matched general-RL reruns.
\end{enumerate}

\section{Recommended language}

A concise methods statement is:

\begin{quote}
We implemented paper-aligned conservation adaptations of PLUS and MOOR under hidden demographics.
PLUS maintained a Bayesian posterior over eight fixed, independently solved Ricker candidate
POMDPs, while MOOR fitted one Ricker model offline to complete ordered episodes and planned in that
same model. Ricker environments provided an in-family control; Allee, theta-logistic, and
regime-switching environments provided preregistered out-of-family misspecification stress tests.
No method received the hidden family or private demographic parameters, and finite-horizon PBVI was
used as a disclosed planner substitution.
\end{quote}

Avoid the phrases ``exact reproduction,'' ``true PLUS,'' or ``true MOOR.'' Use ``paper-aligned
adaptation.'' Retain ``PLUS-inspired'' and ``MOOR-inspired'' for the earlier, easier implementations.

\section{Source and local record}

\begin{itemize}
  \item Memarzadeh and Boettiger (2018), PLUS conservation/fisheries work and archived
  \texttt{pomdpplus} source: \url{https://zenodo.org/records/1161521}.
  \item Ju et al. (2023), ``Model-based offline reinforcement learning for sustainable fishery
  management,'' \emph{Expert Systems}, DOI: \url{https://doi.org/10.1111/exsy.13324}.
  \item Local consolidated record:
  \path{PLUS_MOOR_PAPER_ALIGNMENT_COMPLETE_REFERENCE.md}.
  \item Current reader-facing Markdown:
  \path{PLUS_MOOR_PAPER_ALIGNED_IMPLEMENTATION_EXPLANATION.md}.
  \item Current 72-cell launch design: \path{LAUNCH_PLAN.md}.
  \item Return-blind frozen-runtime acceptance record:
  \path{ACCEPTANCE_V2_FROZEN_RUNTIME_OBSERVABLE_20260719.md}.
\end{itemize}

\section{Conclusion}

The core methodological correction is not merely that the new code uses more complicated equations.
It restores the objects that give PLUS and MOOR their scientific identities. PLUS now reasons with
several fixed, separately solved mechanistic POMDPs and updates a Bayesian posterior over them.
MOOR now fits one mechanistic model through ordered latent trajectories and uses that same model for
prediction, filtering, and planning.

The present design makes the family-level comparison clearer than the earlier cross-family PLUS
extension: both ecological methods assume Ricker, so the three non-Ricker families genuinely test
misspecification. Conservation-specific actions, surveys, rewards, resets, and PBVI are necessary or
disclosed adaptations. They justify the label ``paper-aligned adaptation,'' not exact reproduction.

Implementation verification is now much stronger, but strong empirical conclusions still require a
successful refrozen smoke test, structurally accepted production artifacts, authorized return
analysis, matched general-RL comparisons, and later data/candidate/planner sensitivity studies.

\end{document}
